% ============================================================
% appendix_a.tex — Apéndice A
% ============================================================

\section{Detalles complementarios de configuración}
\label{sec:appendix_a}

\subsection{Ejemplo de prompt de generación con grounding}
\label{subsec:app_prompt}

El siguiente ejemplo ilustra la estructura del prompt utilizado por \sirca{} para la generación contextualizada con grounding documental:

\begin{lstlisting}[basicstyle=\scriptsize\ttfamily, caption={Estructura del prompt de generación con grounding.}, label={lst:prompt}]
[SISTEMA]
Eres un asistente de investigacion biomedica
especializado en plantas medicinales peruanas.
Responde UNICAMENTE basandote en los documentos
proporcionados. Cita cada afirmacion con el
identificador del documento fuente [Doc_ID].
Si no hay informacion suficiente, indicalo
explicitamente.

[CONTEXTO DOCUMENTAL]
[Doc_1] Titulo: Alkaloids from Uncaria tomentosa
Fuente: Journal of Natural Products, 2023
Contenido: "Los alcaloides oxindolicos
pentaciclicos identificados incluyen
mitrafilina, isomitrafilina, pteropodina..."

[Doc_2] Titulo: Anti-inflammatory activity of
cat's claw extracts
Fuente: Phytomedicine, 2022
Contenido: "El extracto estandarizado de
U. tomentosa demostro inhibicion significativa
de NF-kB en modelos in vitro..."

[Doc_3] ...

[CONSULTA DEL USUARIO]
Cuales son los principales alcaloides de
Uncaria tomentosa y que actividad
antiinflamatoria se les atribuye?

[INSTRUCCIONES]
- Fundamenta cada afirmacion citando [Doc_ID]
- No inventes informacion no presente en los
  documentos
- Indica gaps de conocimiento si existen
- Usa terminologia cientifica apropiada
\end{lstlisting}

\subsection{Configuración de conectores de APIs científicas}
\label{subsec:app_apis}

La Tabla~\ref{tab:api_config} presenta la configuración de los conectores utilizados para la adquisición científica automatizada.

\begin{table}[!t]
\centering
\caption{Configuración de conectores de APIs científicas.}
\label{tab:api_config}
\renewcommand{\arraystretch}{1.2}
\footnotesize
\begin{tabular}{l l l}
\toprule
\textbf{API} & \textbf{Endpoint} & \textbf{Rate Limit} \\
\midrule
PubMed E-utilities  & eutils.ncbi.nlm.nih.gov & 3 req/s \\
CrossRef REST       & api.crossref.org         & 50 req/s \\
Semantic Scholar    & api.semanticscholar.org  & 100 req/5min \\
Europe PMC          & europepmc.org/restfulapi & 10 req/s \\
\bottomrule
\end{tabular}
\end{table}

\subsection{Taxonomía de entidades biomédicas reconocidas}
\label{subsec:app_entities}

El módulo de BioNER de \sirca{} reconoce las siguientes categorías de entidades:

\begin{itemize}[leftmargin=*]
  \item \textbf{SPECIES:} Nombres taxonómicos de plantas medicinales (e.g., \textit{Uncaria tomentosa}, \textit{Lepidium meyenii}).
  \item \textbf{CHEMICAL:} Compuestos químicos y metabolitos secundarios (e.g., mitrafilina, macamidas, taspina).
  \item \textbf{ACTIVITY:} Actividades biológicas y farmacológicas (e.g., antiinflamatorio, antioxidante, citotóxico).
  \item \textbf{DISEASE:} Condiciones médicas y patologías (e.g., artritis, cáncer, diabetes).
  \item \textbf{STUDY\_TYPE:} Tipo de estudio (e.g., \textit{in vitro}, \textit{in vivo}, ensayo clínico).
  \item \textbf{TISSUE:} Tejidos y órganos biológicos mencionados en estudios farmacológicos.
\end{itemize}
